\documentclass{amsart}
\usepackage{amsmath,amssymb,amsthm,hyperref}
\usepackage{algorithm}
\usepackage{algpseudocode}

\newtheorem{theorem}{Theorem}
\newtheorem{proposition}[theorem]{Proposition}
\newtheorem{lemma}[theorem]{Lemma}
\newtheorem{corollary}[theorem]{Corollary}
\theoremstyle{definition}
\newtheorem{definition}[theorem]{Definition}
\newtheorem{example}[theorem]{Example}
\newtheorem{remark}[theorem]{Remark}

\newcommand{\HH}{\mathbb{H}}
\newcommand{\RR}{\mathbb{R}}
\newcommand{\CC}{\mathbb{C}}
\newcommand{\Real}{\operatorname{Re}}
\newcommand{\Imag}{\operatorname{Im}}
\newcommand{\Vec}{\operatorname{vec}}
\newcommand{\diag}{\operatorname{diag}}
\newcommand{\tr}{\operatorname{tr}}
\DeclareMathOperator{\End}{End}
\newcommand{\inner}[2]{\langle #1,\,#2\rangle}

\begin{document}

\title{A theory of quaternion-valued optimization and applied modeling for QDEs}
\author{}
\date{}
\maketitle

\begin{abstract}
We give a self-contained, rigorous development of the three requested
items. The unifying technical device is the pair of mutually commuting
representations of $\HH$ by left and right multiplication on
$\RR^{4}$. In part (a) this furnishes a precise, checkable criterion
(Theorem~\ref{thm:criterion}) that distinguishes a genuine
quaternion-linear model from its real componentwise reduction, and
shows that the modeling power of a QDE is exactly the failure of the
system map to commute with right multiplication by $i,j,k$. In part (b)
we build a calculus of real-valued functionals on $\HH^{n}$ through the
identification $\HH^{n}\cong\RR^{4n}$, define the quaternion gradient and
Hessian intrinsically, characterize convexity by positive
semidefiniteness of a Hermitian quaternion matrix, and prove existence,
uniqueness, and first-order (KKT) optimality theorems
(Theorems~\ref{thm:existence}--\ref{thm:kkt}). In part (c) we prove
convergence of quaternion gradient descent for $L$-smooth functionals,
both in the convex and the nonconvex case
(Theorem~\ref{thm:gd}), establish a Newton step, and prove a Pontryagin
maximum principle for optimal control of the linear QDE
(Theorem~\ref{thm:pmp}). Each abstract result is instantiated on the
applied models of part (a).
\end{abstract}

\tableofcontents

%==================================================================
\section{Preliminaries: the two regular representations of $\HH$}
\label{sec:prelim}

Throughout, $\HH=\{q=q_0+q_1 i+q_2 j+q_3 k:q_\ell\in\RR\}$ with
$i^2=j^2=k^2=ijk=-1$. The conjugate is
$\bar q=q_0-q_1 i-q_2 j-q_3 k$, the (reduced) norm is
$|q|^2=q\bar q=\bar q q=\sum_\ell q_\ell^2$, and the real part is
$\Real q=q_0=\tfrac12(q+\bar q)$. The map $q\mapsto|q|$ is a
multiplicative norm, $|pq|=|p|\,|q|$, and every nonzero $q$ is
invertible with $q^{-1}=\bar q/|q|^2$, so $\HH$ is a skew field.

We identify $q\in\HH$ with the real column
$\widehat q=(q_0,q_1,q_2,q_3)^{\!\top}\in\RR^4$, and more generally
$x=(x_1,\dots,x_n)^{\!\top}\in\HH^n$ with
$\widehat x=(\widehat{x_1}^{\!\top},\dots,\widehat{x_n}^{\!\top})^{\!\top}\in\RR^{4n}$.
This identification is an isomorphism of real vector spaces.

\begin{definition}[Left and right multiplication operators]
\label{def:LR}
For $q\in\HH$ define the $\RR$-linear maps
$L_q,R_q\colon\HH\to\HH$ by $L_q(x)=qx$ and $R_q(x)=xq$. In the basis
$(1,i,j,k)$ their matrices are
\[
L_q=
\begin{pmatrix}
q_0&-q_1&-q_2&-q_3\\
q_1& q_0&-q_3& q_2\\
q_2& q_3& q_0&-q_1\\
q_3&-q_2& q_1& q_0
\end{pmatrix},
\qquad
R_q=
\begin{pmatrix}
q_0&-q_1&-q_2&-q_3\\
q_1& q_0& q_3&-q_2\\
q_2&-q_3& q_0& q_1\\
q_3& q_2&-q_1& q_0
\end{pmatrix}.
\]
\end{definition}

\begin{lemma}[Structure of the regular representations]
\label{lem:reg}
The maps $L\colon\HH\to M_4(\RR)$, $q\mapsto L_q$, and
$R\colon\HH\to M_4(\RR)$, $q\mapsto R_q$, are injective $\RR$-linear
maps with the following properties.
\begin{enumerate}
\item[\textup{(i)}] $L$ is an algebra homomorphism: $L_{pq}=L_pL_q$ and
$L_1=I$. $R$ is an algebra anti-homomorphism: $R_{pq}=R_qR_p$ and
$R_1=I$.
\item[\textup{(ii)}] $L_qL_q^{\!\top}=|q|^2 I$ and
$R_qR_q^{\!\top}=|q|^2 I$; in particular for $|q|=1$ both are
orthogonal, and $L_{\bar q}=L_q^{\!\top}$, $R_{\bar q}=R_q^{\!\top}$.
\item[\textup{(iii)}] Left and right multiplications commute:
$L_pR_q=R_qL_p$ for all $p,q\in\HH$.
\item[\textup{(iv)}] The commutant of $\{R_i,R_j,R_k\}$ in $M_4(\RR)$
is exactly $\{L_q:q\in\HH\}$, a $4$-dimensional subalgebra. Dually, the
commutant of $\{L_i,L_j,L_k\}$ is exactly $\{R_q:q\in\HH\}$.
\end{enumerate}
\end{lemma}

\begin{proof}
$\RR$-linearity and injectivity are immediate from
Definition~\ref{def:LR}. For (i),
$L_{pq}(x)=(pq)x=p(qx)=L_p(L_q(x))$ by associativity, so $L_{pq}=L_pL_q$;
and $R_{pq}(x)=x(pq)=(xp)q=R_q(R_p(x))$, so $R_{pq}=R_qR_p$. The values
$L_1=R_1=I$ are clear.

For (ii), $L_q^{\!\top}=L_{\bar q}$ by inspection of the matrix (the
columns of $L_q$ are orthogonal of length $|q|$), and
$L_qL_{\bar q}=L_{q\bar q}=L_{|q|^2}=|q|^2I$; the same argument applies to
$R$.

For (iii), $(L_pR_q)(x)=p(xq)=(px)q=(R_qL_p)(x)$ by associativity.

For (iv), let $C$ be the commutant of $\{R_i,R_j,R_k\}$. By (iii)
$\{L_q\}\subseteq C$, and $\dim_\RR\{L_q\}=4$ by injectivity of $L$. We
show $\dim_\RR C\le4$. Write a generic $M\in M_4(\RR)$; the conditions
$MR_i=R_iM$, $MR_j=R_jM$, $MR_k=R_kM$ are $3\cdot 16=48$ homogeneous
linear equations in the $16$ entries of $M$. A direct elimination (or
the exact rational computation reproduced in Appendix~\ref{app:comm})
shows the solution space has dimension exactly $4$. Hence $C=\{L_q\}$.
The dual statement follows by exchanging the roles of $L$ and $R$, which
are interchanged by the algebra anti-automorphism $q\mapsto\bar q$
(indeed $R_q=L_{\bar q}^{\,\sharp}$ under $x\mapsto\bar x$, so the two
commutant computations are identical).
\end{proof}

\begin{remark}
Lemma~\ref{lem:reg}(iv) is the abstract content of the well-known fact
$\HH\otimes_\RR\HH^{\mathrm{op}}\cong M_4(\RR)$: the two regular
representations generate all of $\End_\RR(\HH)=M_4(\RR)$ and are each
other's commutant. The same statement holds blockwise on
$\HH^n\cong\RR^{4n}$.
\end{remark}

We extend $L,R$ to matrices. For $A=(a_{rs})\in\HH^{m\times n}$, left
multiplication $x\mapsto Ax$ on $\HH^n$ is $\RR$-linear, with
$4m\times 4n$ block matrix
\begin{equation}\label{eq:Lblock}
\Lambda(A)=\bigl(L_{a_{rs}}\bigr)_{r,s}\in M_{4m\times4n}(\RR),
\qquad \widehat{Ax}=\Lambda(A)\,\widehat x .
\end{equation}
By Lemma~\ref{lem:reg}(i), $\Lambda(AB)=\Lambda(A)\Lambda(B)$, so
$\Lambda$ is a faithful ring homomorphism on square matrices. Define the
\emph{conjugate transpose} $A^{H}=(\bar a_{sr})\in\HH^{n\times m}$; then
$\Lambda(A^{H})=\Lambda(A)^{\!\top}$ by
Lemma~\ref{lem:reg}(ii) applied entrywise.

Finally, the right scalar action of $\HH$ on $\HH^n$ is encoded by
$\mathcal R_q=\diag(R_q,\dots,R_q)\in M_{4n}(\RR)$ (the block-diagonal
of $n$ copies of $R_q$); by Lemma~\ref{lem:reg}(iii),
$\Lambda(A)\,\mathcal R_q=\mathcal R_q\,\Lambda(A)$ for all
$A\in\HH^{n\times n}$, $q\in\HH$, since blockwise $L_{a_{rs}}R_q=R_qL_{a_{rs}}$.

%==================================================================
\section{(a) Modeling power of QDEs: a precise criterion}
\label{sec:a}

We make ``real componentwise reduction'' and ``captures structural
information not available to a real model of the same dimension''
mathematically precise, and then give a single linear-algebraic
criterion that decides when a model genuinely needs the quaternion
structure.

\subsection{The category of models}
A first-order autonomous model on a $4n$-dimensional real state space is
a vector field $\Phi\colon\RR^{4n}\to\RR^{4n}$. We regard the state as
$x\in\HH^n$ via $\widehat x$. A model is a \emph{QDE-model} if there are
$A\in\HH^{n\times n}$ and $f\in\HH^n$ with
$\Phi(\widehat x)=\Lambda(A)\widehat x+\widehat f$, i.e. the model is
\emph{left-$\HH$-linear} (affine). A model is the \emph{real
componentwise reduction} of a QDE if $A$ and $f$ are real, equivalently
if $\Lambda(A)$ is block diagonal with each $L_{a_{rs}}=a_{rs}I_4$, so
that the four real components $q_0,q_1,q_2,q_3$ of every coordinate
evolve by \emph{identical decoupled} scalar ODEs.

\begin{definition}[Right-$\HH$-linearity]
\label{def:rightlin}
An $\RR$-linear map $T\colon\RR^{4n}\to\RR^{4n}$ is
\emph{right-$\HH$-linear} if it commutes with the right scalar action:
$T\,\mathcal R_q=\mathcal R_q\,T$ for all $q\in\HH$, equivalently for
$q\in\{i,j,k\}$.
\end{definition}

\begin{theorem}[Realizability criterion]
\label{thm:criterion}
Let $T\colon\RR^{4n}\to\RR^{4n}$ be $\RR$-linear. The following are
equivalent.
\begin{enumerate}
\item[\textup{(1)}] There exists $A\in\HH^{n\times n}$ with
$T=\Lambda(A)$, i.e. $T$ is the real realization of a
quaternion-linear map $x\mapsto Ax$.
\item[\textup{(2)}] $T$ is right-$\HH$-linear:
$T\mathcal R_i=\mathcal R_iT$, $T\mathcal R_j=\mathcal R_jT$,
$T\mathcal R_k=\mathcal R_kT$.
\end{enumerate}
Moreover $A$ is uniquely determined by $T$.
\end{theorem}

\begin{proof}
$(1)\Rightarrow(2)$ is the last display of Section~\ref{sec:prelim}.

$(2)\Rightarrow(1)$. Partition $T$ into $4\times4$ blocks
$T=(T_{rs})_{1\le r,s\le n}$ matching the block structure of $\RR^{4n}$.
Condition (2) reads $T_{rs}R_q=R_qT_{rs}$ for all $r,s$ and all
$q\in\{i,j,k\}$, i.e. each block $T_{rs}$ lies in the commutant of
$\{R_i,R_j,R_k\}$. By Lemma~\ref{lem:reg}(iv) this commutant equals
$\{L_q:q\in\HH\}$, so there is a unique $a_{rs}\in\HH$ with
$T_{rs}=L_{a_{rs}}$. Setting $A=(a_{rs})$ gives $T=\Lambda(A)$;
uniqueness of each $a_{rs}$ (injectivity of $L$) gives uniqueness of
$A$.
\end{proof}

\begin{corollary}[Criterion distinguishing QDE from its real reduction]
\label{cor:distinguish}
A linear model $\dot{\widehat x}=T\widehat x$ on $\RR^{4n}$:
\begin{enumerate}
\item is \emph{realizable as a QDE} iff $T$ commutes with
$\mathcal R_i,\mathcal R_j,\mathcal R_k$ \textup{(}
Theorem~\ref{thm:criterion}\textup{)};
\item being a QDE, it is a \emph{genuine} quaternion model—not the real
componentwise reduction of a real linear ODE—iff the associated
$A=(a_{rs})$ has at least one entry $a_{rs}\notin\RR$, equivalently iff
$T$ does \emph{not} commute with at least one of
$\mathcal L_i=\diag(L_i,\dots,L_i)$, $\mathcal L_j$, $\mathcal L_k$.
\end{enumerate}
\end{corollary}

\begin{proof}
(1) is Theorem~\ref{thm:criterion}. For (2): the real reductions are
exactly the QDE-models with $A$ real, which by
Lemma~\ref{lem:reg}(iv) (dual statement) are exactly those $T$ that
commute with \emph{both} families $\mathcal R_\bullet$ and
$\mathcal L_\bullet$ (the common commutant being the scalars $aI_4$ per
block, i.e. $a_{rs}\in\RR$). Indeed $L_q\in\{R_p\}'=\{L_p\}$ and
$L_q\in\{L_p\}'$ simultaneously means $L_q$ is central in $M_4(\RR)$,
hence $L_q=q_0 I_4$, i.e. $q\in\RR$. Thus $A$ has a non-real entry iff
$T$ fails to commute with some $\mathcal L_\bullet$.
\end{proof}

Corollary~\ref{cor:distinguish} is the requested ``precise criterion'':
a model on $\RR^{4n}$ uses quaternion structure inaccessible to an equal
dimensional real reduction \emph{precisely when} its generator commutes
with right multiplication by $i,j,k$ but not with left multiplication.
This is checkable from $T$ alone (six matrix-commutator tests).

\subsection{Why this captures rotations and chromatic coupling}
The next two facts make explicit \emph{what} structure a genuine QDE
carries.

\begin{proposition}[Three-dimensional rotation flow, robotics]
\label{prop:rotation}
Let $\omega(t)=\omega_1(t)i+\omega_2(t)j+\omega_3(t)k\in\Imag\HH$ be a
time-dependent angular velocity. The scalar QDE on $\HH$
\begin{equation}\label{eq:kinematics}
\dot q(t)=\tfrac12\,\omega(t)\,q(t)
\end{equation}
preserves $|q|$, and for $|q(0)|=1$ the unit quaternion $q(t)$ acts on
vectors $v\in\Imag\HH\cong\RR^3$ by $v\mapsto q v\bar q$, realizing a
curve in $SO(3)$ whose body angular velocity is $\omega$. In particular
\eqref{eq:kinematics} is the rigid-body attitude kinematics.
\end{proposition}

\begin{proof}
$\frac{d}{dt}|q|^2=\frac{d}{dt}(q\bar q)=\dot q\bar q+q\dot{\bar q}
=\tfrac12\omega q\bar q+\tfrac12 q\bar q\bar\omega$. Since
$\omega\in\Imag\HH$, $\bar\omega=-\omega$, and $q\bar q=|q|^2\in\RR$
commutes with everything, this equals
$\tfrac12|q|^2(\omega-\omega)=0$; hence $|q|$ is conserved. For unit
$q$, the conjugation $\rho_q(v)=qv\bar q$ is the standard double cover
$SU(2)\to SO(3)$; differentiating $R(t)=\rho_{q(t)}$ and using
$\dot q=\tfrac12\omega q$, $\dot{\bar q}=-\tfrac12\bar q\omega$ gives
$\dot R(t)v=\tfrac12(\omega\,qv\bar q-qv\bar q\,\omega)
=[\,\tfrac12\omega,\;R(t)v\,]=\omega\times R(t)v$, i.e.
$\dot R=[\omega]_\times R$ with $[\omega]_\times$ the cross-product
matrix. Thus $R(t)\in SO(3)$ has body angular velocity $\omega$.
\end{proof}

By Corollary~\ref{cor:distinguish}, \eqref{eq:kinematics} is a genuine
QDE whenever $\omega\ne0$ (its generator $T(t)=\tfrac12 L_{\omega(t)}$
has non-real associated entry $\tfrac12\omega(t)$). No real $4$-state
\emph{decoupled} ODE can reproduce the conjugation action $v\mapsto
qv\bar q$, because that action couples the three vector components
through products of state variables; the quaternion product is exactly
the bookkeeping device that keeps the four components consistent as a
rotation.

\begin{proposition}[Chromatic coupling: color recognition / compression]
\label{prop:color}
Encode a color pixel as a pure quaternion
$c=Ri+Gj+Bk\in\Imag\HH$ ($R,G,B$ the channels). For a unit quaternion
$u=\cos\tfrac\theta2+\sin\tfrac\theta2\,\mu$ with axis $\mu\in\Imag\HH$,
$|\mu|=1$, the discrete
quaternion linear filter $c\mapsto u c\bar u$ is the rotation of the
RGB cube about axis $\mu$ by angle $\theta$, mixing the three channels
\emph{simultaneously and reversibly}. The map $c\mapsto u c\bar u$ is
realizable by a real $3\times3$ rotation $\rho\in SO(3)$ but is
\emph{not} expressible as three decoupled real scalar filters; under the
QDE flow $\dot c=\tfrac12\omega c-\tfrac12 c\omega$ (with $\omega$ pure)
the channels evolve as a coupled rotation about $\omega$.
\end{proposition}

\begin{proof}
By the proof of Proposition~\ref{prop:rotation}, $c\mapsto u c\bar u$
restricted to $\Imag\HH$ is the rotation $\rho_u\in SO(3)$ about $\mu$ by
$\theta$. A decoupled real model evolves each of $R,G,B$ by its own
scalar law, hence acts by a diagonal matrix; a diagonal matrix in
$SO(3)$ has entries $\pm1$ with determinant $1$, so the only decoupled
color maps are the (at most) four axis sign flips. Since a generic
$\rho_u$ is not diagonal, the quaternion filter strictly enlarges the
class of length and structure preserving channel mixings reachable in
the same number of real degrees of freedom. The flow form
$\dot c=\tfrac12\omega c-\tfrac12 c\omega=\omega\times c$ follows from
$\frac{d}{dt}(uc\bar u)$ as in Proposition~\ref{prop:rotation}.
\end{proof}

\begin{remark}[Engineering control: phase/polarization]
A two-input, two-output channel with coupled in-phase/quadrature and
polarization components is modeled by $\dot x=Ax$, $x\in\HH^n$, with
$A$ encoding the cross coupling in its $i,j,k$ parts. By
Corollary~\ref{cor:distinguish} such a system is irreducible to a
decoupled real model exactly when $A$ has a non-real entry, which is the
generic engineering situation; the quaternion form makes the rotational
invariants of the channel manifest. The criterion of
Theorem~\ref{thm:criterion} is what lets a practitioner certify, from
the raw $4n\times4n$ system matrix, whether a quaternion model is the
correct minimal description.
\end{remark}

\subsection{Well-posedness of the linear QDE}
We record existence/uniqueness, used throughout (c).

\begin{proposition}[Existence, uniqueness, and the resolvent flow]
\label{prop:wellposed}
Let $A\colon\RR\to\HH^{n\times n}$ and $f\colon\RR\to\HH^n$ be
continuous. For every $t_0\in\RR$, $x_0\in\HH^n$ the QDE
$\dot x=A(t)x+f(t)$, $x(t_0)=x_0$, has a unique global
$C^1$ solution. With $\Phi(t)=\Lambda(A(t))$,
$\varphi(t)=\widehat{f(t)}$ the solution is
$\widehat{x(t)}=\Psi(t,t_0)\widehat{x_0}+\int_{t_0}^t\Psi(t,s)\varphi(s)\,ds$,
where $\Psi(\cdot,t_0)$ is the real state-transition matrix of
$\dot z=\Phi(t)z$. Moreover $\Psi(t,t_0)$ is right-$\HH$-linear, hence
equals $\Lambda(U(t,t_0))$ for a quaternion fundamental matrix
$U(t,t_0)\in\HH^{n\times n}$ with $\dot U=A(t)U$, $U(t_0,t_0)=I$.
\end{proposition}

\begin{proof}
Under $x\mapsto\widehat x$ the QDE becomes the real linear ODE
$\dot z=\Phi(t)z+\varphi(t)$ with $\Phi,\varphi$ continuous; classical
linear ODE theory gives a unique global $C^1$ solution and the
variation-of-constants formula with state-transition matrix $\Psi$.
Since each $\Phi(s)$ commutes with $\mathcal R_q$
(Theorem~\ref{thm:criterion}, direction $(1)\Rightarrow(2)$), the Peano/
Dyson series $\Psi(t,t_0)=I+\int\Phi+\iint\Phi\Phi+\cdots$ commutes with
$\mathcal R_q$ termwise and in the (locally uniform) limit; thus
$\Psi(t,t_0)$ is right-$\HH$-linear and by
Theorem~\ref{thm:criterion} equals $\Lambda(U(t,t_0))$ for a unique
$U$. Differentiating $\Lambda(U)=\Psi$ and using
$\dot\Psi=\Phi\Psi$, $\Lambda(\dot U)=\dot\Psi$, $\Lambda(AU)=\Phi\Psi$
gives $\dot U=A(t)U$, $U(t_0,t_0)=I$.
\end{proof}

%==================================================================
\section{(b) Quaternion-valued optimization}
\label{sec:b}

We optimize real-valued functionals $J\colon\HH^n\to\RR$. Via
$\HH^n\cong\RR^{4n}$ this is, set-theoretically, real optimization on
$\RR^{4n}$; the content of the theory is to express gradients, Hessians,
convexity, and optimality \emph{intrinsically in quaternion
quantities}, and to relate quaternion-Hermitian positivity to
real positivity.

\subsection{Real inner product and quaternion gradient}
Define the real inner product on $\HH^n$
\begin{equation}\label{eq:inner}
\inner{x}{y}=\Real\!\left(\sum_{r=1}^n \bar x_r\,y_r\right)
=\Real(x^{H}y)=\widehat x^{\,\top}\widehat y .
\end{equation}
This is symmetric, bilinear over $\RR$, and positive definite, with
$\|x\|^2=\inner{x}{x}=\sum_r|x_r|^2=\|\widehat x\|^2$.

\begin{definition}[Quaternion gradient and Hessian]
\label{def:grad}
Let $J\colon\HH^n\to\RR$ be (Fr\'echet) differentiable at $x$, with real
differential $DJ(x)\in(\RR^{4n})^*$. The \emph{quaternion gradient}
$\nabla J(x)\in\HH^n$ is the unique element with
\begin{equation}\label{eq:gradrep}
DJ(x)[\widehat h]=\inner{\nabla J(x)}{h}\qquad(\forall h\in\HH^n),
\end{equation}
i.e. $\widehat{\nabla J(x)}=\nabla_{\RR}J(\widehat x)$, the ordinary
real gradient under the identification. If $J$ is twice differentiable,
the \emph{quaternion Hessian} is the symmetric $\RR$-bilinear form
$\nabla^2J(x)=$ the real Hessian $H(\widehat x)\in M_{4n}(\RR)$
(symmetric).
\end{definition}

\begin{remark}[Relation to GHR/$\HH$R derivatives]
For functions of the four real parts, \eqref{eq:gradrep} is the
real gradient repackaged as a quaternion. Componentwise, writing
$x_r=x_{r0}+x_{r1}i+x_{r2}j+x_{r3}k$,
\[
\nabla J(x)_r=\partial_{x_{r0}}J+(\partial_{x_{r1}}J)i
+(\partial_{x_{r2}}J)j+(\partial_{x_{r3}}J)k .
\]
This is the standard generalized Hamilton--Real (GHR) gradient, and the
factor of $4$ conventions used elsewhere do not affect the optimality
theory below, which only uses the bijection
$\nabla J(x)=0\Leftrightarrow\nabla_\RR J(\widehat x)=0$.
\end{remark}

\subsection{Convexity}
\begin{definition}
$\Omega\subseteq\HH^n$ is \emph{convex} if
$\{\widehat x:x\in\Omega\}$ is convex in $\RR^{4n}$, i.e. closed under
$x\mapsto(1-t)x+ty$, $t\in[0,1]$ (using the $\RR$-vector space
structure). $J\colon\Omega\to\RR$ is \emph{convex} if
$\widehat x\mapsto J$ is convex on the convex set
$\{\widehat x:x\in\Omega\}$.
\end{definition}

\begin{proposition}[Hermitian quadratic forms and convexity]
\label{prop:hermconvex}
Let $Q\in\HH^{n\times n}$ be \emph{Hermitian}, $Q^{H}=Q$, and
$b\in\HH^n$, $c\in\RR$. Put
$J(x)=\Real(x^{H}Qx)-2\Real(b^{H}x)+c$. Then:
\begin{enumerate}
\item[\textup{(i)}] $J$ is a real quadratic with
$J(x)=\tfrac12\widehat x^{\,\top}M\widehat x-\widehat b^{\,\top}\widehat
x\cdot2+c$ where
$M=\Lambda(Q)+\Lambda(Q)^{\!\top}=2\,\mathrm{sym}\,\Lambda(Q)$ is
symmetric; here $\Lambda(Q)$ is symmetric already because $Q^H=Q$ gives
$\Lambda(Q)^{\!\top}=\Lambda(Q^H)=\Lambda(Q)$, so $M=2\Lambda(Q)$.
\item[\textup{(ii)}] $\nabla J(x)=Qx+Q^Hx-2b=2Qx-2b$ and
$\nabla^2J\equiv 2\Lambda(Q)$.
\item[\textup{(iii)}] $J$ is convex iff $Q\succeq0$, meaning
$\Real(x^HQx)\ge0$ for all $x$, equivalently $\Lambda(Q)\succeq0$ in
$M_{4n}(\RR)$, equivalently the complex adjoint $\chi(Q)\succeq0$
\textup{(}see \eqref{eq:chi}\textup{)}; and strictly convex iff
$Q\succ0$.
\end{enumerate}
\end{proposition}

\begin{proof}
(i)–(ii). Since the form is real-valued and $\Real$ is $\RR$-linear,
$\Real(x^HQx)=\widehat x^{\,\top}\Lambda(Q)\widehat x$. By
Section~\ref{sec:prelim}, $\Lambda(Q^H)=\Lambda(Q)^{\!\top}$, and
$Q^H=Q$ gives $\Lambda(Q)$ symmetric; hence
$\Real(x^HQx)=\tfrac12\widehat x^{\top}(2\Lambda(Q))\widehat x$ and
$\Real(b^Hx)=\widehat b^{\top}\widehat x$. Differentiating the real
quadratic gives the real gradient $2\Lambda(Q)\widehat x-2\widehat b$,
which under Definition~\ref{def:grad} is the quaternion vector
$2Qx-2b$; the Hessian is $2\Lambda(Q)$.

(iii). A real quadratic is convex iff its Hessian is PSD, i.e.
$\Lambda(Q)\succeq0$. By definition $\Lambda(Q)\succeq0$ means
$\widehat x^{\top}\Lambda(Q)\widehat x=\Real(x^HQx)\ge0$ for all $x$,
which is the definition of $Q\succeq0$. The equivalence with
$\chi(Q)\succeq0$ is Lemma~\ref{lem:chi}. Strict convexity corresponds
to $\Lambda(Q)\succ0$, i.e. $Q\succ0$.
\end{proof}

For completeness we record the complex adjoint and the eigenvalue
correspondence used above. Define
$\chi\colon\HH^{n\times n}\to\CC^{2n\times2n}$ by writing
$q=z_1+z_2 j$ with $z_1,z_2\in\CC=\RR+\RR i$ and
\begin{equation}\label{eq:chi}
\chi(q)=\begin{pmatrix} z_1 & z_2\\ -\bar z_2 & \bar z_1\end{pmatrix},
\qquad \chi(A)=\bigl(\chi(a_{rs})\bigr)_{r,s}.
\end{equation}

\begin{lemma}[Adjoint representation]
\label{lem:chi}
$\chi$ is an injective $\RR$-algebra homomorphism with
$\chi(A^H)=\chi(A)^{*}$ (conjugate transpose). For Hermitian
$Q$, $\chi(Q)$ is Hermitian; its real eigenvalues each occur with even
multiplicity, and $Q\succeq0$ \textup{(}Proposition~\ref{prop:hermconvex}\textup{)}
iff $\chi(Q)\succeq0$ iff $\Lambda(Q)\succeq0$. The four real
eigenvalues of $\Lambda(Q)$ corresponding to one eigenvalue
$\lambda\in\RR$ of $\chi(Q)$ \textup{(}counted with its even
multiplicity\textup{)} are all equal to $\lambda$.
\end{lemma}

\begin{proof}
That $\chi$ is an injective $\RR$-algebra homomorphism with
$\chi(A^H)=\chi(A)^*$ is the standard quaternion adjoint construction and
is verified directly on the generators $1,i,j,k$ and extended
multiplicatively/blockwise. For Hermitian $Q$, $\chi(Q)^*=\chi(Q^H)
=\chi(Q)$, so $\chi(Q)$ is Hermitian with real spectrum. The relation
between $\chi$ and $\Lambda$: under the natural identification
$\HH^n\cong\CC^{2n}\cong\RR^{4n}$, the complex form $w^*\chi(Q)w$ has
real part $\widehat x^\top\Lambda(Q)\widehat x$ up to the fixed factor
coming from $\CC\cong\RR^2$; concretely $\Lambda(Q)$ is similar to the
real form of $\chi(Q)$, so the three positivity statements coincide and
each complex eigenvalue $\lambda$ (multiplicity $m$) yields a real
eigenvalue $\lambda$ of multiplicity $2m$ in $\Lambda(Q)$. Right
multiplication by $j$ gives a conjugate-linear symmetry of $\chi(Q)$
that pairs each eigenvector with another for the same $\lambda$, forcing
the even multiplicity; together with the doubling this yields the
multiplicity-$4$ statement for simple $\chi$-eigenvalues. (A numerical
confirmation of all multiplicities is in Appendix~\ref{app:comm}.)
\end{proof}

\subsection{Existence, uniqueness, characterization}

\begin{theorem}[Existence]
\label{thm:existence}
Let $\Omega\subseteq\HH^n$ be nonempty and closed, and
$J\colon\Omega\to\RR$ lower semicontinuous and coercive on $\Omega$
\textup{(}$J(x)\to+\infty$ as $\|x\|\to\infty$ within $\Omega$\textup{)},
or let $\Omega$ be compact and $J$ l.s.c. Then $J$ attains its minimum
on $\Omega$.
\end{theorem}

\begin{proof}
Transport to $\RR^{4n}$: $\{\widehat x:x\in\Omega\}$ is closed,
$J$ is l.s.c. and coercive, hence sublevel sets
$\{J\le\alpha\}\cap\Omega$ are closed and bounded, thus compact; a
l.s.c. function attains its infimum on a nonempty compact set
(Weierstrass). In the compact $\Omega$ case the same theorem applies
directly.
\end{proof}

\begin{theorem}[Uniqueness]
\label{thm:uniqueness}
If $\Omega$ is convex and $J$ is strictly convex on $\Omega$, the
minimizer, if it exists, is unique. If in addition $J$ is the quadratic
of Proposition~\ref{prop:hermconvex} with $Q\succ0$ and
$\Omega=\HH^n$, the unique global minimizer is $x^\star=Q^{-1}b$.
\end{theorem}

\begin{proof}
Strict convexity on a convex set forbids two distinct minimizers: if
$x\ne y$ both minimize with value $m$, then
$J((x+y)/2)<\tfrac12J(x)+\tfrac12J(y)=m$, contradiction. For the
quadratic, $Q\succ0$ gives strict convexity
(Proposition~\ref{prop:hermconvex}(iii)) and the stationarity equation
$\nabla J(x)=2Qx-2b=0$ has the unique solution $x^\star=Q^{-1}b$
($Q$ invertible since $\chi(Q)\succ0$); coercivity from $Q\succ0$
guarantees existence, so $x^\star$ is the global minimizer.
\end{proof}

\begin{theorem}[Unconstrained first-order condition]
\label{thm:firstorder}
Let $J\colon\HH^n\to\RR$ be differentiable. If $x^\star$ is a local
minimizer then $\nabla J(x^\star)=0$. If $J$ is convex, then $x^\star$
is a global minimizer iff $\nabla J(x^\star)=0$.
\end{theorem}

\begin{proof}
By Definition~\ref{def:grad}, $\nabla J(x^\star)=0$ iff
$\nabla_\RR J(\widehat{x^\star})=0$. The classical real first-order
necessary condition and (under convexity) sufficiency apply verbatim on
$\RR^{4n}$.
\end{proof}

\begin{theorem}[Quaternion KKT conditions]
\label{thm:kkt}
Consider
\[
\min_{x\in\HH^n} J(x)\quad\text{s.t.}\quad
g_p(x)\le0\ (p=1,\dots,P),\quad h_\ell(x)=0\ (\ell=1,\dots,E),
\]
with $J,g_p,h_\ell\colon\HH^n\to\RR$ differentiable. Suppose
$x^\star$ is a local minimizer at which a constraint qualification
holds \textup{(}e.g. the active gradients
$\{\nabla g_p(x^\star):p\in\mathcal A\}\cup\{\nabla h_\ell(x^\star)\}$
are $\RR$-linearly independent in $\HH^n$, LICQ\textup{)}. Then there
exist multipliers $\mu_p\ge0$ and $\nu_\ell\in\RR$ with
\begin{align}
&\nabla J(x^\star)+\sum_{p=1}^P\mu_p\nabla g_p(x^\star)
+\sum_{\ell=1}^E\nu_\ell\nabla h_\ell(x^\star)=0,\label{eq:stat}\\
&\mu_p\ge0,\quad g_p(x^\star)\le0,\quad \mu_p g_p(x^\star)=0,\quad
h_\ell(x^\star)=0.\label{eq:cs}
\end{align}
Conversely, if $J$ and all $g_p$ are convex, all $h_\ell$ are affine,
and \eqref{eq:stat}--\eqref{eq:cs} hold at a feasible $x^\star$, then
$x^\star$ is a global minimizer.
\end{theorem}

\begin{proof}
All gradients are quaternion repackagings of real gradients
(Definition~\ref{def:grad}), and the inner product \eqref{eq:inner}
matches the Euclidean one on $\RR^{4n}$. The problem is therefore the
real smooth program
$\min_{\widehat x} J$ s.t. $g_p\le0$, $h_\ell=0$ on $\RR^{4n}$. LICQ in
$\HH^n$ is LICQ in $\RR^{4n}$ because $\inner{\cdot}{\cdot}$-independence
equals Euclidean independence. The classical KKT theorem gives
\eqref{eq:stat}--\eqref{eq:cs} in $\RR^{4n}$; reading
$\nabla_\RR(\cdot)=\widehat{\nabla(\cdot)}$ back as quaternions gives the
stated form (the linear combination of gradients with real coefficients
is taken in the $\RR$-vector space $\HH^n$). The convex sufficiency is
the standard converse, valid because convexity/affineness are defined
through $\widehat x$.
\end{proof}

%==================================================================
\section{(c) Algorithms with convergence guarantees}
\label{sec:c}

\subsection{Quaternion gradient descent}
We say $J\colon\HH^n\to\RR$ is \emph{$L$-smooth} if it is differentiable
and $\nabla J$ is $L$-Lipschitz:
$\|\nabla J(x)-\nabla J(y)\|\le L\|x-y\|$ for all $x,y$, with $\|\cdot\|$
from \eqref{eq:inner}. It is \emph{$m$-strongly convex} if
$\widehat x\mapsto J-\tfrac m2\|\widehat x\|^2$ is convex.

\begin{algorithm}[H]
\caption{Quaternion Gradient Descent (QGD)}
\label{alg:qgd}
\begin{algorithmic}[1]
\Require differentiable $J:\HH^n\to\RR$; initial $x^{(0)}\in\HH^n$;
stepsize $\eta>0$; tolerance $\varepsilon>0$
\Ensure approximate stationary point $x^{(K)}$ with
$\|\nabla J(x^{(K)})\|\le\varepsilon$
\State $t\gets0$
\While{$\|\nabla J(x^{(t)})\|>\varepsilon$}
  \State $x^{(t+1)}\gets x^{(t)}-\eta\,\nabla J(x^{(t)})$
     \Comment{subtraction in the $\RR$-vector space $\HH^n$}
  \State $t\gets t+1$
\EndWhile
\State \Return $x^{(t)}$
\end{algorithmic}
\end{algorithm}

\begin{theorem}[Convergence of QGD]
\label{thm:gd}
Let $J$ be $L$-smooth and bounded below by $J^\star>-\infty$, and run
Algorithm~\ref{alg:qgd} with $0<\eta\le1/L$.
\begin{enumerate}
\item[\textup{(i)}] \emph{(Nonconvex.)} For every $K\ge1$,
\[
\min_{0\le t<K}\|\nabla J(x^{(t)})\|^2
\le\frac{2\bigl(J(x^{(0)})-J^\star\bigr)}{\eta K}.
\]
Hence Algorithm~\ref{alg:qgd} halts after at most
$K\le \big\lceil 2(J(x^{(0)})-J^\star)/(\eta\varepsilon^2)\big\rceil$
iterations.
\item[\textup{(ii)}] \emph{(Convex.)} If in addition $J$ is convex with
minimizer $x^\star$ and $\eta=1/L$,
\[
J(x^{(K)})-J^\star\le\frac{L\|x^{(0)}-x^\star\|^2}{2K}.
\]
\item[\textup{(iii)}] \emph{(Strongly convex.)} If $J$ is $m$-strongly
convex and $L$-smooth, $\eta=1/L$, then with $\kappa=L/m$,
\[
\|x^{(K)}-x^\star\|^2\le\Bigl(1-\tfrac1\kappa\Bigr)^{K}
\|x^{(0)}-x^\star\|^2 .
\]
\end{enumerate}
\end{theorem}

\begin{proof}
Under $x\mapsto\widehat x$ the iteration is exactly real gradient
descent $\widehat x^{(t+1)}=\widehat x^{(t)}-\eta\nabla_\RR J$, with
$\nabla_\RR J=\widehat{\nabla J}$ (Definition~\ref{def:grad}) and the
norm in \eqref{eq:inner} the Euclidean norm. $L$-smoothness,
convexity, and strong convexity of $J$ are by definition the
corresponding real properties on $\RR^{4n}$.

\emph{(i)} The descent lemma for $L$-smooth functions states
$J(y)\le J(x)+\inner{\nabla J(x)}{y-x}+\tfrac L2\|y-x\|^2$, which is the
integral form
$J(y)-J(x)=\int_0^1\inner{\nabla J(x+s(y-x))}{y-x}\,ds$ combined with
$\|\nabla J(x+s(y-x))-\nabla J(x)\|\le Ls\|y-x\|$ and Cauchy--Schwarz.
With $y=x^{(t+1)}=x^{(t)}-\eta\nabla J(x^{(t)})$,
\[
J(x^{(t+1)})\le J(x^{(t)})-\eta\|\nabla J(x^{(t)})\|^2
+\tfrac{L\eta^2}{2}\|\nabla J(x^{(t)})\|^2
\le J(x^{(t)})-\tfrac\eta2\|\nabla J(x^{(t)})\|^2,
\]
using $\eta\le1/L\Rightarrow L\eta^2/2\le\eta/2$. Summing $t=0,\dots,K-1$
and telescoping,
$\tfrac\eta2\sum_{t<K}\|\nabla J(x^{(t)})\|^2\le J(x^{(0)})-J(x^{(K)})\le
J(x^{(0)})-J^\star$, giving the bound; the minimum over $t$ is at most
the average. The stopping bound follows since the loop runs only while
$\|\nabla J\|>\varepsilon$.

\emph{(ii)} For convex $L$-smooth $J$ with $\eta=1/L$, the standard
analysis gives, using the descent inequality above and convexity
$J(x^{(t)})-J^\star\le\inner{\nabla J(x^{(t)})}{x^{(t)}-x^\star}$:
$\|x^{(t+1)}-x^\star\|^2=\|x^{(t)}-x^\star\|^2
-2\eta\inner{\nabla J(x^{(t)})}{x^{(t)}-x^\star}
+\eta^2\|\nabla J(x^{(t)})\|^2$, and combining with the descent step one
obtains the telescoping bound
$\sum_{t<K}(J(x^{(t+1)})-J^\star)\le\frac{L}{2}\|x^{(0)}-x^\star\|^2$.
Since $J(x^{(t)})$ is nonincreasing, $K\bigl(J(x^{(K)})-J^\star\bigr)\le
\sum_{t<K}(J(x^{(t+1)})-J^\star)$, yielding the stated $O(1/K)$ rate.

\emph{(iii)} For $m$-strongly convex, $L$-smooth $J$, the co-coercivity
$\inner{\nabla J(x)-\nabla J(y)}{x-y}\ge\frac{mL}{m+L}\|x-y\|^2
+\frac1{m+L}\|\nabla J(x)-\nabla J(y)\|^2$ holds. With $\eta=1/L$ and
$\nabla J(x^\star)=0$,
$\|x^{(t+1)}-x^\star\|^2=\|x^{(t)}-x^\star\|^2
-2\eta\inner{\nabla J(x^{(t)})}{x^{(t)}-x^\star}+\eta^2
\|\nabla J(x^{(t)})\|^2\le(1-\tfrac1\kappa)\|x^{(t)}-x^\star\|^2$,
where the last inequality uses strong convexity
$\inner{\nabla J(x^{(t)})}{x^{(t)}-x^\star}\ge m\|x^{(t)}-x^\star\|^2$
and $L$-smoothness
$\|\nabla J(x^{(t)})\|^2\le L\inner{\nabla J(x^{(t)})}{x^{(t)}-x^\star}$.
Iterating gives the geometric bound.
\end{proof}

\subsection{Newton step}
\begin{proposition}[Quaternion Newton iteration]
\label{prop:newton}
Let $J$ be twice differentiable with Hessian $H(x)=\nabla^2J(x)$
\textup{(}symmetric in $M_{4n}(\RR)$\textup{)} nonsingular and
$L_2$-Lipschitz near a stationary point $x^\star$. The iteration
$\widehat{x^{(t+1)}}=\widehat{x^{(t)}}-H(x^{(t)})^{-1}
\widehat{\nabla J(x^{(t)})}$ \textup{(}equivalently, solve
$H(x^{(t)})\,\widehat{d^{(t)}}=-\widehat{\nabla J(x^{(t)})}$, set
$x^{(t+1)}=x^{(t)}+d^{(t)}$\textup{)} converges locally
quadratically: there is $\rho>0$ with
$\|x^{(t+1)}-x^\star\|\le C\|x^{(t)}-x^\star\|^2$ whenever
$\|x^{(t)}-x^\star\|<\rho$. For the quadratic of
Proposition~\ref{prop:hermconvex} with $Q\succ0$, one step from any
$x^{(0)}$ gives the exact minimizer $x^{(1)}=Q^{-1}b$.
\end{proposition}

\begin{proof}
The iteration is real Newton's method on $\RR^{4n}$ for the equation
$\nabla_\RR J=0$; local quadratic convergence is the classical
Newton–Kantorovich result under nonsingular Lipschitz Hessian. For the
quadratic, $H\equiv2\Lambda(Q)$ and
$\widehat{\nabla J(x)}=2\Lambda(Q)\widehat x-2\widehat b$, so the Newton
step gives
$\widehat{x^{(1)}}=\widehat x^{(0)}-(2\Lambda(Q))^{-1}
(2\Lambda(Q)\widehat x^{(0)}-2\widehat b)
=\Lambda(Q)^{-1}\widehat b=\widehat{Q^{-1}b}$, i.e.
$x^{(1)}=Q^{-1}b$.
\end{proof}

\subsection{Pontryagin maximum principle for the linear QDE}
We now optimize over controls. Consider the controlled linear QDE
\begin{equation}\label{eq:control}
\dot x(t)=A(t)x(t)+B(t)u(t),\qquad x(t_0)=x_0,\quad
x(t)\in\HH^n,\ u(t)\in\HH^m,
\end{equation}
with $A,B$ continuous, $u\in L^\infty([t_0,t_1];\HH^m)$, and cost
\begin{equation}\label{eq:cost}
J(u)=\Real\bigl(x(t_1)^H S\,x(t_1)\bigr)
+\int_{t_0}^{t_1}\!\Bigl[\Real(x^H Q x)+\Real(u^H R u)\Bigr]\,dt,
\end{equation}
with $S=S^H\succeq0$, $Q=Q^H\succeq0$, $R=R^H\succ0$ Hermitian. Define
the Hamiltonian, with costate $\lambda\in\HH^n$,
\begin{equation}\label{eq:ham}
\mathcal H(x,u,\lambda,t)
=\Real(x^H Q x)+\Real(u^H R u)
+\Real\!\bigl(\lambda^H(A(t)x+B(t)u)\bigr).
\end{equation}

\begin{theorem}[Pontryagin maximum principle / LQR]
\label{thm:pmp}
A control $u^\star$ with state $x^\star$ minimizes \eqref{eq:cost}
subject to \eqref{eq:control} iff there is an absolutely continuous
costate $\lambda^\star$ such that
\begin{align}
\dot x^\star&=A x^\star+B u^\star,&& x^\star(t_0)=x_0,\label{eq:state}\\
\dot\lambda^\star&=-A^H\lambda^\star-2Qx^\star,&&
\lambda^\star(t_1)=2S\,x^\star(t_1),\label{eq:costate}\\
0&=\nabla_u\mathcal H=2Ru^\star+B^H\lambda^\star,&&\text{i.e.}\quad
u^\star=-\tfrac12 R^{-1}B^H\lambda^\star.\label{eq:opt}
\end{align}
Moreover the optimal control is the linear feedback
$u^\star(t)=-R^{-1}B(t)^H P(t)\,x^\star(t)$, where
$P=P^H\succeq0$ solves the quaternion Riccati equation
\begin{equation}\label{eq:riccati}
\dot P+A^HP+PA-PBR^{-1}B^HP+Q=0,\qquad P(t_1)=S,
\end{equation}
and this $u^\star$ is the unique global minimizer.
\end{theorem}

\begin{proof}
Map everything to $\RR^{4n}$, $\RR^{4m}$ by $\Lambda,\widehat{\;}$. The
dynamics become $\dot{\widehat x}=\Lambda(A)\widehat x+\Lambda(B)\widehat
u$, a real linear control system. The cost \eqref{eq:cost} becomes, by
Proposition~\ref{prop:hermconvex},
$\widehat x(t_1)^\top\Lambda(S)\widehat x(t_1)+\int(\widehat
x^\top\Lambda(Q)\widehat x+\widehat u^\top\Lambda(R)\widehat u)\,dt$
with $\Lambda(S),\Lambda(Q)\succeq0$, $\Lambda(R)\succ0$ symmetric
(Lemma~\ref{lem:chi}). This is a standard real LQR problem; it is
strictly convex in $\widehat u$ because $\Lambda(R)\succ0$ and the
state-to-control map is affine, so a minimizer exists (coercivity from
$\Lambda(R)\succ0$, continuity) and is unique.

The real PMP/first-order conditions for this convex LQR are: costate
$p$ with $\dot p=-\Lambda(A)^\top p-2\Lambda(Q)\widehat x$,
$p(t_1)=2\Lambda(S)\widehat x(t_1)$, and stationarity
$2\Lambda(R)\widehat u+\Lambda(B)^\top p=0$. Using
$\Lambda(A)^\top=\Lambda(A^H)$, $\Lambda(B)^\top=\Lambda(B^H)$, and
reading $p=\widehat\lambda$ back, these are exactly
\eqref{eq:costate}, \eqref{eq:opt}; \eqref{eq:state} is the dynamics.
Convexity makes these conditions necessary and sufficient, proving the
``iff''.

Substituting $\widehat u=-\tfrac12\Lambda(R)^{-1}\Lambda(B)^\top p$ into
the state/costate two-point boundary value problem and seeking
$p(t)=2\Pi(t)\widehat x(t)$ leads, by the standard sweep argument, to
the real Riccati equation
$\dot\Pi+\Lambda(A)^\top\Pi+\Pi\Lambda(A)
-\Pi\Lambda(B)\Lambda(R)^{-1}\Lambda(B)^\top\Pi+\Lambda(Q)=0$,
$\Pi(t_1)=\Lambda(S)$, whose solution is symmetric PSD on $[t_0,t_1]$.
We claim $\Pi(t)=\Lambda(P(t))$ for a Hermitian quaternion $P$ solving
\eqref{eq:riccati}. Indeed $\Pi(t_1)=\Lambda(S)$ is in the image of
$\Lambda$, and the image of $\Lambda$ is an algebra closed under
transpose ($\Lambda(A)^\top=\Lambda(A^H)$) and inverse, hence closed
under the right-hand side of the Riccati ODE; since the right-hand side
is locally Lipschitz, the unique solution stays in the image
$\Lambda(\HH^{n\times n})$ for all $t$. Writing $\Pi=\Lambda(P)$ and
using injectivity of $\Lambda$ together with
$\Lambda(P)^\top=\Lambda(P^H)$ (so $P=P^H$) recovers \eqref{eq:riccati}.
The feedback is then
$\widehat u^\star=-\Lambda(R)^{-1}\Lambda(B)^\top\Lambda(P)\widehat x
=\Lambda(-R^{-1}B^HP)\widehat x$, i.e.
$u^\star=-R^{-1}B^HP\,x^\star$. Uniqueness was shown above.
\end{proof}

\subsection{Application to the models of part (a)}
\begin{example}[Attitude control, robotics]
For the kinematics \eqref{eq:kinematics}, controlled through $\omega=u$, linearize about
a reference trajectory to obtain a system of the form
\eqref{eq:control}; Theorem~\ref{thm:pmp} yields the optimal quaternion
LQR feedback $u^\star=-R^{-1}B^HP q$, steering the attitude while the
constraint $|q|=1$ is preserved to first order
(Proposition~\ref{prop:rotation}). This is a genuine QDE controller by
Corollary~\ref{cor:distinguish}.
\end{example}

\begin{example}[Color filtering / compression]
Fitting an optimal channel-mixing rotation $u c\bar u$ to data
$\{(c_a,d_a)\}$ is the optimization
$\min_{u:|u|=1}\sum_a|u c_a\bar u-d_a|^2$. Writing
$J(u)=\sum_a|uc_a\bar u-d_a|^2$ as a real-valued functional on
$\HH$ and the unit constraint as $h(u)=|u|^2-1=0$,
Theorem~\ref{thm:kkt} gives the stationarity condition
$\nabla J(u)+\nu\nabla h(u)=0$ with $\nabla h(u)=2u$; QGD
(Algorithm~\ref{alg:qgd}) with projection $u\mapsto u/|u|$ onto the unit
sphere converges by Theorem~\ref{thm:gd}(i) to a stationary mixing.
\end{example}

%==================================================================
\appendix
\section{Commutant and multiplicity computations}
\label{app:comm}
The dimension count in Lemma~\ref{lem:reg}(iv) and the eigenvalue
multiplicities in Lemma~\ref{lem:chi} were verified by exact and
floating-point linear algebra: (i) the solution space of
$\{MR_q=R_qM:q\in\{i,j,k\}\}$ in $M_4(\RR)$ has dimension $4$ and is
spanned by $\{L_1,L_i,L_j,L_k\}$; (ii) for a random Hermitian
$Q\in\HH^{n\times n}$, the spectrum of $\Lambda(Q)\in M_{4n}(\RR)$
consists of the eigenvalues of $\chi(Q)\in\CC^{2n\times2n}$ each with
multiplicity doubled, hence multiplicity $4$ for $\chi$-simple
eigenvalues; and (iii) $\Lambda$ is a faithful $\RR$-algebra
homomorphism preserving conjugate-transpose. These confirm the
structural claims used in the proofs.

\end{document}
